
\newgeometry{left=3cm,right=3cm,top=2.5cm,bottom=2cm,includefoot}
\fontsize{10pt}{12pt}\selectfont
\begin{abstract}
    \large{\textbf{FAST/FIX протоколын судалгаа}} \par
        Хөрөнгийн зах зээлд мэдээллийн урсгалын хурд, найдвартай байдал нь арилжааны гүйцэтгэл, эрсдэлийн удирдлага, үйл ажиллагааны үр ашигт шууд нөлөөлдөг. Иймээс олон улсын жишигт нийцсэн мессежийн стандартыг судалж, байгууллагын дотоод процессийг автоматжуулах нь зүй ёсны хэрэгцээ болж байна. Энэхүү дадлага ажлын хүрээнд зах зээлийн мэдээлэл ба арилжааны системүүдийн хооронд өгөгдөл солилцох түгээмэл стандарт болох FIX (Financial Information eXchange) болон дамжуулалтын үр ашгийг нэмэгдүүлэхэд зориулагдсан FAST (FIX Adapted for STreaming) аргачлалыг судалсан. Хоёр протокол нь XML хэл дээр бичигддэг. Эдгээр протоколыг ашиглан Монголын хөрөнгийн бирж(Цаашид, МХБ) үнэт цаасны арилжааг гүйцэтгэдэг.\par
  \textbf{- Үндэслэл} \par
    БиДиСек компанийн хувьд одоогийн хэрэглэдэг системээ өөр гуравдагч талын гүйцэтгэгч компаниар гүйцэтгүүлсэн. Энэ нь хуучирсан, эх код өөрсдөд нь байхгүй нь улмаас үргэлж бусдаас хараат байдгаас үүдэн системээ шинэчлэх хэрэгтэй болсон.
    \par
  \textbf{- Зорилго}
  \begin{itemize}
      \item FIX болон FAST-ийн архитектур, мессежийн бүтэц, кодчилол/декодчиллын зарчим, гүйцэтгэлийн шинжүүдийг системтэйгээр нэгтгэн ойлгох.
      \item Хэрэглээний түгээмэл кейсүүдийг тодорхойлж, хэрэгжүүлэх боломжит загвар, сайн туршлагын зөвлөмж боловсруулах.
  \end{itemize}
  
  \textbf{- Зорилт}
    \begin{itemize}
        \item Албан стандарт, техникийн тайлбар бичгүүдийг судалж, ойлголтын зураглал гаргах.
        \item FIX-ийн сэдвүүд (сесс, апликейшн давхарга, мессеж төрлүүд, tag=value формат) ба FAST шахалтын механизмыг харьцуулан тайлбарлах.
        \item Саатал, зурвас өргөн, найдвартай байдал, аюулгүй байдлын үзүүлэлтүүдийн онцлогийг нэгтгэх.
        \item Жишээ мессежүүд дээр кодчилол/декодчиллын туршилт хийх (нээлттэй сан ашиглах боломжтой бол).
    \end{itemize}

   \large{\textbf{n8n автоматжуулалт}} \par
   \textbf{Үндэслэл} \par
   \begin{itemize}
       \item Гараар давтагддаг ажлыг бууруулж, өгөгдлийн урсгалыг найдвартай, давтамжтайгаар гүйцэтгэх хэрэгцээ өндөр.
       \item Код бичилт бага шийдлээр (low-code) олон эх үүсвэрийг холбож, логжилт, мэдэгдэл, алдааны барьцалдалтыг стандартчлах шаардлагатай.
       \item Өөрчлөлт, өргөтгөл хийхэд хурдан, засвар үйлчилгээ хялбар архитектур хэрэгтэй.
   \end{itemize}
   \textbf{Зорилго} \par
   \begin{itemize}
       \item n8n дээр байгууллагын хэрэгцээнд нийцсэн автоматжуулалтын урсгалуудыг загварчилж, хэрэгжүүлэх.
        \item Интеграцийн зангилаа, алдааны менежмент, мониторингийг тогтворжуулсан дахин ашиглагдах шийдэл болгох.
   \end{itemize}

    \textbf{Зорилт} \par
    \begin{itemize}
       \item Шаардлага цуглуулж, хэрэглэх зангилаа (Webhook, API, DB, Google Sheets, Redis, мэдэгдэл г.м)-гийн зураглалыг гаргах.
       \item Өгөгдөл хувиргалт, баталгаажуулалт, нөхцөлт салбарлалт (Switch-IF), дахин оролдлого, тайм-аутын дүрмийг тогтоох.
       \item Туршилтын урсгалууд бүтээж, төгсгөл хүртэлх гүйцэтгэлийг лог/экзекюшн тайлангаар баталгаажуулах.
       \item Мониторинг ба алдааны мэдэгдлийн сувгийг (имэйл/чат/вэбхүүк) холбож, ажиллагааны гажилтыг хурдан илрүүлэх.
   \end{itemize}


\end{abstract}